\clearpage\section{MOS transistor schematic}
These editable vector schematics connect equal node names within a circuit. Primitive-cell names are local pins; the hierarchy defines their connections. B denotes the body terminal; 0 is global ground. Each panel identifies its netlist and scope; measurement sources are connected as described.
Exact decoder primitives and the scalable hierarchy of \texttt{01\_two\_to\_four\_load.sp}, \texttt{02\_three\_to\_eight\_truth.sp}, \texttt{03\_eight\_to\_256\_truth.sp} and \texttt{04\_eight\_to\_256\_timing.sp}.
\par
\begin{figure}[!ht]
\centering
\begin{circuitikz}[font=\small]
\node[pmos,bulk] (MP0) at (0.0,3) {};
\node[anchor=west,font=\scriptsize] at ($(MP0.center)+(.55,.40)$) {MP0};
\draw (MP0.bulk) -- ++(.30,0) node[right,font=\scriptsize] {B=vp};
\draw (MP0.G) -- ++(-.45,0) node[left,font=\small] {\texttt{a0}};
\draw (MP0.S) -- (MP0.S |- 0,4.2); \draw (MP0.D) -- (MP0.D |- 0,1.7);
\node[nmos,bulk] (MN0) at (1.9,0.0) {};
\node[anchor=west,font=\scriptsize] at ($(MN0.center)+(.55,.40)$) {MN0};
\draw (MN0.bulk) -- ++(.30,0) node[right,font=\scriptsize] {B=vn};
\draw (MN0.G) -- ++(-.45,0) node[left,font=\small] {\texttt{a0}};
\node[pmos,bulk] (MP1) at (3.8,3) {};
\node[anchor=west,font=\scriptsize] at ($(MP1.center)+(.55,.40)$) {MP1};
\draw (MP1.bulk) -- ++(.30,0) node[right,font=\scriptsize] {B=vp};
\draw (MP1.G) -- ++(-.45,0) node[left,font=\small] {\texttt{a1}};
\draw (MP1.S) -- (MP1.S |- 0,4.2); \draw (MP1.D) -- (MP1.D |- 0,1.7);
\node[nmos,bulk] (MN1) at (1.9,-2.2) {};
\node[anchor=west,font=\scriptsize] at ($(MN1.center)+(.55,.40)$) {MN1};
\draw (MN1.bulk) -- ++(.30,0) node[right,font=\scriptsize] {B=vn};
\draw (MN1.G) -- ++(-.45,0) node[left,font=\small] {\texttt{a1}};
\draw (MN0.S) -- (MN1.D) node[midway,above right,font=\small] {\texttt{n1}};
\draw (MP0.S |- 0,4.2) -- (MP1.S |- 0,4.2) node[midway,above]{vp}; \draw (MP0.D |- 0,1.7) -- (MP1.D |- 0,1.7); \draw (MN0.D) -- (MN0.D |- 0,1.7) node[circ,label=above right:z]{};
\draw (MN1.S) -- ++(0,-0.4) node[below,font=\small] {\texttt{vn}};
\end{circuitikz}
\caption{Exact NAND2 cell. AND2 connects NAND2.z to inv.a; its output is inv.z.}
\end{figure}
\par
\par
\begin{figure}[!ht]
\centering
\begin{circuitikz}[font=\small]
\node[pmos,bulk] (MP) at (0,2.5) {};
\node[anchor=west,font=\scriptsize] at ($(MP.center)+(.55,.40)$) {MP};
\draw (MP.bulk) -- ++(.30,0) node[right,font=\scriptsize] {B=vp};
\node[nmos,bulk] (MN) at (0,0.2) {};
\node[anchor=west,font=\scriptsize] at ($(MN.center)+(.55,.40)$) {MN};
\draw (MN.bulk) -- ++(.30,0) node[right,font=\scriptsize] {B=vn};
\draw (MP.G) -- (-1,2.5) -- (-1,.2) -- (MN.G); \draw (-1,1.35) node[circ]{} -- (-1.6,1.35) node[left]{a}; \draw (MP.D |- 0,1.35) node[circ]{} -- (2,1.35) node[right]{z};
\draw (MP.S) -- ++(0,0.5) node[above,font=\small] {\texttt{vp}};
\draw (MP.D) -- (MN.D) ;
\draw (MN.S) -- ++(0,-0.5) node[below,font=\small] {\texttt{vn}};
\end{circuitikz}
\caption{Exact inv cell used for input complements and NAND-output restoration.}
\end{figure}
\par
All $L=0.25\,\mu$m. inv uses $W_P/W_N=3.6k/1.8k\,\mu$m; nand2 uses $3.6k/3.6k\,\mu$m. PMOS sources/bodies connect to vp, NMOS bodies to vn. Named n1 is the internal NMOS stack node, not a supply rail.
\clearpage
\par
\begin{figure}[!ht]
\centering
\begin{circuitikz}[font=\small]
\node[pmos,bulk] (MP1) at (0,3) {};
\node[anchor=west,font=\scriptsize] at ($(MP1.center)+(.55,.40)$) {MP1};
\draw (MP1.bulk) -- ++(.30,0) node[right,font=\scriptsize] {B=vp};
\draw (MP1.G) -- ++(-.45,0) node[left,font=\small] {\texttt{a1}};
\node[pmos,bulk] (MP0) at (0,0.6) {};
\node[anchor=west,font=\scriptsize] at ($(MP0.center)+(.55,.40)$) {MP0};
\draw (MP0.bulk) -- ++(.30,0) node[right,font=\scriptsize] {B=vp};
\draw (MP0.G) -- ++(-.45,0) node[left,font=\small] {\texttt{a0}};
\draw (MP1.S) -- ++(0,0.4) node[above,font=\small] {\texttt{vp}};
\draw (MP1.D) -- (MP0.S) node[midway,above right,font=\small] {\texttt{p1}};
\draw (MP0.D) -- ++(0,-0.4) node[below,font=\small] {\texttt{z}};
\node[nmos,bulk] (MN0) at (4.5,1.8) {};
\node[anchor=west,font=\scriptsize] at ($(MN0.center)+(.55,.40)$) {MN0};
\draw (MN0.bulk) -- ++(.30,0) node[right,font=\scriptsize] {B=vn};
\draw (MN0.G) -- ++(-.45,0) node[left,font=\small] {\texttt{a0}};
\draw (MN0.D) -- ++(0,0.4) node[above,font=\small] {\texttt{z}};
\draw (MN0.S) -- ++(0,-0.4) node[below,font=\small] {\texttt{vn}};
\node[nmos,bulk] (MN1) at (8.0,1.8) {};
\node[anchor=west,font=\scriptsize] at ($(MN1.center)+(.55,.40)$) {MN1};
\draw (MN1.bulk) -- ++(.30,0) node[right,font=\scriptsize] {B=vn};
\draw (MN1.G) -- ++(-.45,0) node[left,font=\small] {\texttt{a1}};
\draw (MN1.D) -- ++(0,0.4) node[above,font=\small] {\texttt{z}};
\draw (MN1.S) -- ++(0,-0.4) node[below,font=\small] {\texttt{vn}};
\end{circuitikz}
\caption{Exact NOR2 alternative: MP1 lies above MP0 in the PMOS stack.}
\end{figure}
\par
NOR2 uses $W_P=7.2k\,\mu$m and $W_N=1.8k\,\mu$m. MP0's source is p1, while its body remains vp. For the 2-to-4 NAND version, each positive minterm is NAND2 followed by inv. With input order $(a1,a0)$, the NAND input pairs for outputs 0--3 are $(nb1,nb0)$, $(nb1,a0)$, $(a1,nb0)$ and $(a1,a0)$. The NOR alternative uses $(a1,a0)$, $(a1,rb0)$, $(rb1,a0)$ and $(rb1,rb0)$ respectively. nb/rb are generated by the architecture's own two input inverters; vp maps to supply vn or vr and the cell's vn maps to global 0. Both variants attach cl from every output to 0.
\par
\begin{figure}[!ht]
\centering
\begin{circuitikz}[font=\small]
\node[draw,rounded corners=2pt,align=center,minimum width=2.4cm,minimum height=.9cm] (I) at (0,0) {$a[2:0]$ + inv\\3 complements};
\node[draw,rounded corners=2pt,align=center,minimum width=2.5cm,minimum height=.9cm] (N) at (4,0) {8 NAND3\\minterms};
\node[draw,rounded corners=2pt,align=center,minimum width=2.5cm,minimum height=.9cm] (O) at (8,0) {8 inv\\one-hot outputs};
\draw[-{Latex}] (I.east) -- (N.west);
\draw[-{Latex}] (N.east) -- (O.west);
\end{circuitikz}
\caption{3-to-8 decoder: all eight minterms are explicit NAND3 + inv cells.}
\end{figure}
\par
nand3 extends the shown NAND2 pattern with three parallel PMOS and three series NMOS; $W_P/W_N=3.6k/5.4k\,\mu$m. All k=1; 3 input inverters plus 8 AND3 cells total 70 MOS, each output loaded by 10 fF. The full 8-to-256 topology is shown on the next page without drawing 256 duplicate output cells.
\clearpage
\par
\begin{figure}[!ht]
\centering
\begin{circuitikz}[font=\small]
\node[draw,rounded corners=2pt,align=center,minimum width=3cm,minimum height=.9cm] (A) at (0,0) {$a[7:0]$\\8 inv: tb0--tb7};
\node[draw,rounded corners=2pt,align=center,minimum width=3.4cm,minimum height=.9cm] (B) at (4.6,0) {2 groups of 16 AND4\\tp0 / tp1\\$k=2$};
\node[draw,rounded corners=2pt,align=center,minimum width=2.8cm,minimum height=.9cm] (C) at (9.2,0) {256 AND2\\t0--t255\\$k=1$};
\draw[-{Latex}] (A.east) -- (B.west);
\draw[-{Latex}] (B.east) -- (C.west);
\end{circuitikz}
\caption{Two logical stages after input complement generation. Each AND is NAND + inv.}
\end{figure}
\par
\par
\begin{figure}[!ht]
\centering
\begin{circuitikz}[font=\small]
\node[draw,rounded corners=2pt,align=center,minimum width=3cm,minimum height=.9cm] (A) at (0,0) {$a[7:0]$\\8 inv: hb0--hb7};
\node[draw,rounded corners=2pt,align=center,minimum width=3.4cm,minimum height=.9cm] (B) at (4.6,0) {4 groups of 4 AND2\\hg0--hg3\\$k=1$};
\draw[-{Latex}] (A.east) -- (B.west);
\node[draw,rounded corners=2pt,align=center,minimum width=3.4cm,minimum height=.9cm] (C) at (4.6,-2.5) {2 groups of 16 AND2\\hp0 / hp1\\$k=2$};
\node[draw,rounded corners=2pt,align=center,minimum width=2.8cm,minimum height=.9cm] (D) at (9.2,-2.5) {256 AND2\\h0--h255\\$k=1$};
\draw[-{Latex}] (B.south) -- (C.north);
\draw[-{Latex}] (C.east) -- (D.west);
\end{circuitikz}
\caption{Three logical stages, with exact grouping and sizing multipliers.}
\end{figure}
\par
For integer output index j, the two-stage path combines tp0[$j\bmod16$] (bits 3:0) and tp1[$\lfloor j/16\rfloor$] (bits 7:4). The three-stage path first decodes bit pairs (1:0), (3:2), (5:4), (7:6), combines adjacent pairs into hp0/hp1, then uses the same low/high nibble indices. Every AND2 or AND4 is the exact NAND-plus-inverter combination. nand4 uses four parallel PMOS ($3.6k\,\mu$m) and four series NMOS ($7.2k\,\mu$m), all $L=0.25\,\mu$m; k also scales the following inverter. PMOS body rails map to vt for the two-stage design and vh for the three-stage design; all NMOS bodies connect to 0. Every output has cl=10 fF at baseline. Counts including eight input inverters are 1872 and 1840 MOS. The 2/3 logical stages contain 4/6 electrical gate levels after the input literals; complemented-input paths add one level, giving 5/7. Supply/input sources and repeated capacitors are omitted from the hierarchy boxes but retained in the exact source decks.
